\section{Three orthogonal design axes}
\label{sec:axes}

The cleanest way to read the comparison is that each system makes an
independent choice on three axes, and each axis decides exactly one of
the outcomes a user observes. \Cref{tab:axes} is the whole comparison in
one table; the rest of the note justifies it.

\begin{table}[t]
\centering
\renewcommand{\arraystretch}{1.25}
\begin{tabularx}{\textwidth}{@{}lXXl@{}}
\toprule
\textbf{Axis} & \textbf{\binius{}} & \textbf{\zinc{}} & \textbf{Decides} \\
\midrule
Arithmetization &
  unrolled word \emph{circuit} (arbitrary wiring) &
  uniform \emph{AIR} over $\Ftwo[X]$ (structured wiring) &
  \textbf{verifier} \\
Commitment / proximity &
  Reed--Solomon $+$ recursive FRI/BaseFold &
  linear $\Ftwo$-RAA $+$ flat Brakedown &
  \textbf{proof size} \\
Reduction field &
  sumcheck tensors in $\mathrm{GF}(2^{128})$ &
  discharge kept in $\Ftwo/\mathrm{GF}(2^8)$ &
  \textbf{memory / large $N$} \\
\bottomrule
\end{tabularx}
\caption{The three design axes. Each system picks a different point on
each axis, and each axis decides one headline outcome. No single axis
decides the prover --- it is dominated by content both systems share (the
$\AND$ discharge and the commitment encode); at the matched rate the net
of offsetting per-phase gaps leaves \zinc{} a touch ahead and widening
with $N$.}
\label{tab:axes}
\end{table}

\subsection{Arithmetization decides the verifier}
\label{sec:axis-arith}

A verifier is succinct exactly when the constraint system's structure ---
its ``wiring,'' the map from each constraint to the witness positions it
reads --- is \emph{succinctly evaluable}: computable at a random point in
$\polylog(N)$ without a commitment. This is a property of the
arithmetization, not the commitment.

\begin{itemize}
  \item A \emph{uniform AIR} has structured wiring: the same transition
  relation at every row, with inter-row references that are fixed shifts
  of the trace domain and intra-word references that are fixed monomials.
  The verifier evaluates one fixed relation at a random point --- an
  $O(\log N)$ formula --- and never reads the per-row structure.

  \item A \emph{word circuit} has arbitrary wiring: each gate may draw any
  committed word, shifted any way. There is no formula for ``which word
  feeds gate $j$''; it is data --- the circuit description. To evaluate
  the wiring at a random point the verifier sums over all $N$ gates, in
  the clear. This is \binius{}'s $O(N)$ verifier
  (\Cref{sec:shifts,sec:results} make it concrete and measured).
\end{itemize}

SHA-256 happens to be uniform (64 identical rounds plus a uniform message
schedule), which is why \zinc{} can use an AIR and why prime-field STARKs
do too. \binius{} pays an $O(N)$ verifier as the deliberate price of a
\emph{general} circuit model that also handles non-uniform computations.
We return in \Cref{sec:analysis} to whether \binius{} could take an AIR on
the SHA path specifically, and what it would cost.

The same axis decides \emph{preprocessing}. A word circuit's wiring is
data, so its constraint-system and key description grow $O(N)$ --- the
hundreds of megabytes \Cref{sec:res-scaling} records, and the cause of
\binius{}'s circuit-build wall. A uniform AIR's structure is a fixed
formula, so \zinc{}'s preprocessing is $O(1)$ (a UAIR specification plus
the code seeds, kilobyte-scale and independent of the compression count).
The $O(N)$-versus-$O(1)$ preprocessing gap and the linear-versus-$\sqrt N$
verifier are two faces of the same arithmetization choice.

\subsection{Commitment decides proof size}
\label{sec:axis-commit}

The proximity test is independent of the arithmetization. \zinc{} buys a
\emph{linear-time} encoder (an $\Ftwo$-RAA code, no NTT) and pays for it
with a \emph{flat} proof: a Brakedown-style test opens $t \approx 987$
full columns, so the proof grows as $t\cdot\sqrt{N}$ with a large
$t$-floor. \binius{} buys a \emph{small, recursive} proof (FRI folds the
Reed--Solomon codeword in $\log N$ halvings, a few hashes per round) and
pays for it with an NTT encoder. The proof-size axis is therefore the
flat-versus-recursive proximity choice, and it is the gap that
\Cref{sec:results} shows \binius{} winning, widening with $N$. The fix on
the \zinc{} side is a WHIR-style recursive proximity test
\cite{whir}, independent of everything else.

\subsection{Reduction field decides memory and scaling}
\label{sec:axis-field}

The third axis is the field in which the nonlinearity is reduced.
\binius{}'s BitAnd and Shift reductions build $O(N)$ tensors in
$\mathrm{GF}(2^{128})$; together with a rate-$\tfrac12$ Reed--Solomon
codeword, the working set spills cache super-linearly at large $N$.
\zinc{} keeps the witness bit-packed and runs the Hadamard discharge in the
base field $\Ftwo/\mathrm{GF}(2^8)$, lifting to $\mathrm{GF}(2^{128})$
only at the $\psi$ projections --- roughly $16$--$32\times$ fewer bytes
streamed. The consequence is memory: at the largest matched instance
\zinc{} uses about half the footprint, and its prover scales sub-linearly
where \binius{}'s goes super-linear on a memory-constrained host
(\Cref{sec:results}).

\subsection{Why no single axis decides the prover}

No single axis decides the prover, because the prover is dominated by
work both systems must do: discharge the same $\AND$ content and encode
the witness for commitment. The word-level (one packed lane) versus
bit-sliced ($\times\wbit$) asymmetry in how the two reach that content
roughly cancels at the current optimisation level, so the prover is a
\emph{net} of large offsetting per-phase gaps rather than a single
dominant difference. At the matched commitment rate $1/4$ that net leaves
\zinc{} a touch ahead at moderate $N$; because \zinc{} scales linearly and
\binius{} super-linearly, the gap widens to $\sim$$2\times$ at scale. The
per-phase decomposition in \Cref{sec:results} shows the offsetting gaps
explicitly, and \Cref{sec:analysis} traces how they net out with $N$.
